\documentclass[11pt]{article}

\usepackage[margin=1.05in]{geometry}
\usepackage[T1]{fontenc}
\usepackage[utf8]{inputenc}
\usepackage{lmodern}
\usepackage{microtype}
\usepackage{amsmath,amssymb,amsthm,mathtools}
\usepackage{bm}
\usepackage{enumitem}
\usepackage{xcolor}
\usepackage[colorlinks=true,linkcolor=blue!45!black,citecolor=blue!45!black,urlcolor=blue!45!black]{hyperref}

\newtheorem{proposition}{Proposition}
\newtheorem{remark}{Remark}
\newcommand{\dd}{\,\mathrm{d}}
\newcommand{\ii}{\mathrm{i}}
\newcommand{\TT}{\mathbb{T}}
\newcommand{\CC}{\mathbb{C}}
\newcommand{\BZ}{\mathrm{BZ}}

\title{Floquet Nielsen--Ninomiya on an Extended Torus:\\A Cleaned Working Note}
\author{P. Shubham Parashar}
\date{Cleaned note, June 2026}

\begin{document}
\maketitle

\begin{abstract}
This note is a cleaned, conservative version of an exploratory note on Nielsen--Ninomiya-type constraints for Floquet or spacetime-lattice fermions.  The central example is the scalar kernel
\[
  D(\omega,k)=\mu-e^{\ii(\omega-k)}
\]
on the extended torus \(\TT^2=S^1_\omega\times S^1_k\).  The useful calculation is an elementary loop winding of \(D\) along the non-contractible cycle \(\theta=\omega-k\).  The note also states carefully what this calculation does and does not prove: it is not a counterexample to the static Hermitian Nielsen--Ninomiya theorem, but a compact way of seeing how the relevant object changes when one works with a frequency-momentum/Floquet kernel rather than a static Hermitian Hamiltonian \(K(k)\) on the spatial Brillouin zone.
\end{abstract}

\section{Purpose and cleanup verdict}

The original note mixed three related but logically distinct ideas:
\begin{enumerate}[label=(\roman*)]
  \item Friedan's formulation of the Nielsen--Ninomiya theorem for a static Hermitian lattice Hamiltonian;
  \item a scalar Floquet/spacetime kernel \(D(\omega,k)=\mu-e^{\ii(\omega-k)}\);
  \item the claim that a nonzero winding of \(D\) should be read as a chiral index.
\end{enumerate}

The cleaned version keeps the explicit calculation but weakens the interpretation.  The robust statement is:

\begin{quote}
The scalar kernel \(D(\omega,k)=\mu-e^{\ii(\omega-k)}\) defines a map from \(\TT^2\) to \(\CC^\times\) for \(|\mu|\neq1\).  Along a one-cycle on which \(\theta=\omega-k\) winds once, the phase of \(D\) has winding number \(+1\) for \(|\mu|<1\) and \(0\) for \(|\mu|>1\).
\end{quote}

This is best treated as a technical working note, not as a publication claim.  In particular, the invariant here is a \emph{one-cycle winding}; it is not a two-dimensional Chern number on \(\TT^2\), and by itself it does not constitute a theorem-level evasion of Nielsen--Ninomiya.

\section{Static Nielsen--Ninomiya: the object being constrained}

A standard version of the Nielsen--Ninomiya theorem concerns a static, translationally invariant, local, Hermitian lattice Hamiltonian
\begin{equation}
  H=\int_{\BZ}\dd^d p\; \psi_a^\dagger(\bm p)K^a{}_b(\bm p)\psi^b(\bm p),
\end{equation}
where \(K(\bm p)=K^\dagger(\bm p)\) is a smooth matrix-valued function on the spatial Brillouin torus \(\TT^d\).  Near an isolated Fermi point \(\bm p_\alpha\), one assumes a relativistic block
\begin{equation}
  K(\bm p)\simeq (p-p_\alpha)^i\Gamma_i^\alpha
\end{equation}
up to gapped spectator bands.  The chirality of the low-energy fermion is determined by the orientation of this Dirac cone, or equivalently by the sign of the appropriate product of gamma matrices.

Friedan's formulation packages the statement in terms of a Green's-function current in \((p^0,\bm p)\)-space.  The important structural assumptions are:
\begin{enumerate}[label=(\alph*)]
  \item locality and smoothness of \(K(\bm p)\) on the compact spatial torus;
  \item Hermiticity of the static Hamiltonian kernel;
  \item charge conservation and block decomposition into irreducible charge sectors;
  \item relativistic low-energy behavior near isolated zeros.
\end{enumerate}
Under these assumptions the total chirality in each charge sector vanishes.  Slogan: a static local Hermitian lattice Hamiltonian on the spatial Brillouin zone cannot produce an unpaired chiral fermion.

\section{The Floquet/spacetime toy kernel}

Now change the object.  Instead of a static Hamiltonian \(K(k)\) on a spatial circle \(S^1_k\), consider a scalar kernel on the extended torus \(S^1_\omega\times S^1_k\),
\begin{equation}
  D(\omega,k)=\mu-e^{\ii(\omega-k)},
  \label{eq:D}
\end{equation}
where \(\omega\) is a dimensionless frequency or quasienergy angle and \(k\) is crystal momentum.  This is not a Hermitian Hamiltonian kernel.  It is better thought of as a frequency-momentum Dirac operator or inverse propagator.

Equation \eqref{eq:D} depends only on
\begin{equation}
  \theta=\omega-k \mod 2\pi.
\end{equation}
Therefore the topology of the example is the topology of a map from a one-cycle into \(\CC^\times\), even though the kernel is written on \(\TT^2\).

\section{Correct winding computation}

Restrict to a loop \(C\subset\TT^2\) on which \(\theta=\omega-k\) winds once.  Define
\begin{equation}
  f(\theta)=\mu-e^{\ii\theta},
  \qquad \theta\in[0,2\pi].
\end{equation}
For \(|\mu|\neq1\), this loop does not pass through the origin in the complex plane.  The winding number is
\begin{equation}
  \nu(C)=\frac{1}{2\pi\ii}\int_C \dd\log D
       =\frac{1}{2\pi\ii}\int_0^{2\pi}\dd\theta\,\frac{f'(\theta)}{f(\theta)}.
  \label{eq:winding}
\end{equation}
With \(z=e^{\ii\theta}\), one has \(f=\mu-z\), \(\dd\theta=\dd z/(\ii z)\), and
\begin{equation}
  \frac{f'(\theta)}{f(\theta)}\dd\theta
  =\frac{-\ii z}{\mu-z}\frac{\dd z}{\ii z}
  =\frac{\dd z}{z-\mu}.
\end{equation}
Therefore
\begin{equation}
  \nu(C)=\frac{1}{2\pi\ii}\oint_{|z|=1}\frac{\dd z}{z-\mu}
  =
  \begin{cases}
    1, & |\mu|<1,\\[2pt]
    0, & |\mu|>1.
  \end{cases}
\end{equation}

\begin{proposition}
For the scalar kernel \(D(\omega,k)=\mu-e^{\ii(\omega-k)}\), the pullback of \(D\) to any loop on which \(\theta=\omega-k\) winds once has winding number \(+1\) for \(|\mu|<1\) and \(0\) for \(|\mu|>1\).  Reversing the orientation of the loop reverses the sign.
\end{proposition}

\begin{remark}
The sign is cycle-dependent.  If \(\omega\) winds once with \(k\) fixed, then \(\theta\) winds once and the invariant is \(+1\).  If \(k\) winds once with \(\omega\) fixed, then \(\theta\) winds in the opposite direction and the invariant is \(-1\) in the \(|\mu|<1\) phase.
\end{remark}

\section{What this does and does not show}

The one-form
\begin{equation}
  a=\frac{1}{2\pi\ii}\dd\log D
\end{equation}
is closed wherever \(D\neq0\).  Its periods over one-cycles of the torus are integers.  This is the clean mathematical content of the example.

It is tempting to write a Friedan-like current
\begin{equation}
  j^\mu=\frac{1}{2\pi\ii}\epsilon^{\mu\nu}\partial_\nu\log D.
\end{equation}
This is useful as notation, but it should not be oversold.  For \(|\mu|\neq1\), the kernel has no zero on the real torus.  Thus there is no literal two-dimensional divergence source sitting on \(\TT^2\).  The nontrivial invariant is instead the period of \(\dd\log D\) around a non-contractible one-cycle.

This distinction matters.  Since \(\CC^\times\simeq S^1\), one has \(\pi_1(\CC^\times)=\mathbb Z\), which supports loop windings, but \(\pi_2(\CC^\times)=0\), so there is no scalar two-dimensional winding number of a map \(\TT^2\to\CC^\times\) analogous to a Chern number.

\section{Relation to chiral propagation}

The kernel \eqref{eq:D} also has a simple Fourier-series interpretation.  For \(|\mu|<1\),
\begin{equation}
  \frac{1}{\mu-e^{\ii\theta}}
  =-e^{-\ii\theta}\sum_{n\ge0}\mu^n e^{-\ii n\theta}.
\end{equation}
For \(|\mu|>1\),
\begin{equation}
  \frac{1}{\mu-e^{\ii\theta}}
  =\frac{1}{\mu}\sum_{n\ge0}\mu^{-n}e^{\ii n\theta}.
\end{equation}
Thus the two regimes correspond to opposite choices of allowed Fourier support along the discrete diagonal direction associated with \(\theta=\omega-k\).  With a fixed convention for Fourier transforms, this gives propagation along one directed lattice ray.  This is the sense in which the toy kernel resembles a chiral spacetime-lattice propagator.

The continuum statement \(G(x,t)\propto\theta(t)\delta(x-vt)\) should be treated as a heuristic continuum limit, not as an exact result of the compact-torus model without specifying boundary conditions, Fourier conventions, and regularization.

\section{Comparison with Nielsen--Ninomiya hypotheses}

The safest way to state the relation to Nielsen--Ninomiya is a hypotheses table:

\begin{center}
\begin{tabular}{p{0.42\linewidth}p{0.42\linewidth}}
\hline
Static Nielsen--Ninomiya object & Floquet/spacetime toy object \\
\hline
Hermitian \(K(k)\) on spatial \(S^1_k\) & Non-Hermitian/inverse-propagator-like \(D(\omega,k)\) on \(S^1_\omega\times S^1_k\) \\
Local static Hamiltonian & Frequency-momentum kernel / spacetime-lattice kernel \\
Chirality from zeros of \(K(k)\) and Dirac-cone orientation & Loop winding of \(D\) around \(\CC^\times\) \\
No-go theorem constrains net static chirality on the spatial BZ & The toy winding lies outside the static Hermitian hypotheses \\
\hline
\end{tabular}
\end{center}

Therefore the note should not say that the toy model ``violates'' Nielsen--Ninomiya.  It should say:

\begin{quote}
The toy model lies outside the assumptions of the static Hermitian Nielsen--Ninomiya theorem.  Its nontrivial integer is a loop winding of a frequency-momentum kernel on an extended torus, not the net chirality of a static Hamiltonian on the spatial Brillouin zone.
\end{quote}

\section{Possible upgrades}

To turn this into a stronger project, one would need at least one of the following:
\begin{enumerate}[label=(\alph*)]
  \item a fully specified unitary Floquet operator \(U(k)\) whose inverse propagator reduces to \eqref{eq:D};
  \item a precise Ginsparg--Wilson or overlap-fermion dictionary relating \(D\) to a lattice chiral symmetry;
  \item a higher-dimensional analogue where the invariant is a genuine Chern number or Green's-function index;
  \item a clean anomaly calculation showing quantized charge pumping or spectral flow.
\end{enumerate}

Until then, this note is best used as a compact worked example and a reminder of the difference between static Hamiltonian topology and extended frequency-momentum topology.

\begin{thebibliography}{9}

\bibitem{NielsenNinomiya}
H. B. Nielsen and M. Ninomiya,
\emph{Absence of Neutrinos on a Lattice. II. Intuitive Topological Proof},
Nucl. Phys. B \textbf{193}, 173--194 (1981).

\bibitem{Friedan}
D. Friedan,
\emph{A Proof of the Nielsen--Ninomiya Theorem},
Commun. Math. Phys. \textbf{85}, 481--490 (1982).

\bibitem{DeMarcoWen}
M. DeMarco and X.-G. Wen,
\emph{A Single Right-Moving Free Fermion Mode on an Ultra-Local 1+1D Spacetime Lattice with Exactly Lorentz Invariant Dynamics},
\href{https://arxiv.org/abs/1805.03663}{arXiv:1805.03663}.

\end{thebibliography}

\end{document}
